% Physical unit 013; accepted Session 02.
\GordanUnitStart{013}{printed page 8}
\GordanTopFolio{8}
\begin{center}
{\large CHAPTER III.}\\[0.8em]
\S\ 1.
\end{center}

Let there be named:
\[
 a\int_0^{\frac\pi2} d\varphi\sqrt{1-e^2+e^2\sin^2 l_0\cos^2\varphi}
\]
the period of extension, and
\[
 \int_0^{\frac\pi2}\cos l_0\,d\varphi\,
 {\sqrt{1-e^2+e^2\sin^2 l_0\cos^2\varphi}\over 1-\sin^2 l_0\cos^2\varphi}
\]
the period of longitude of the geodetic line $l_0$; since these depend only on the angle $l_0$, let them be denoted by $\Sigma(l_0)$ and $\mathcal A(l_0)$. Since the values of these periods will hereafter be of very great importance, the series must now be developed by which they are computed easily and with the greatest possible accuracy.

\begin{center}\S\ 2.\end{center}

For finding $\Sigma(l_0)$ we expand under the sign of integration
\[
 \sqrt{1-e^2+e^2\sin^2 l_0\cos^2\varphi}
\]
according to the binomial theorem, where, for brevity, we put:
\[
 {1\over4}\,{e^2\over1-e^2}={\varepsilon\over(1-\varepsilon)^2}=o.
\]
It then follows:
\[
\begin{aligned}
 {1\over b}\Sigma(l_0)
 &=\int_0^{\frac\pi2}d\varphi\{1+2o\sin^2l_0\cos^2\varphi-2o^2\sin^4l_0\cos^4\varphi\\
 &\qquad{}+4o^3\sin^6l_0\cos^6\varphi-10o^4\sin^8l_0\cos^8\varphi\ldots\}\\
 &=\int_0^{\frac\pi2}d\varphi+2o\sin^2l_0\int_0^{\frac\pi2}\cos^2\varphi\,d\varphi\\
 &\quad{}-2o^2\sin^4l_0\int_0^{\frac\pi2}\cos^4\varphi\,d\varphi
 +4o^3\sin^6l_0\int_0^{\frac\pi2}\cos^6\varphi\,d\varphi\\
 &\quad{}-10o^4\sin^8l_0\int_0^{\frac\pi2}\cos^8\varphi\,d\varphi\ldots
\end{aligned}
\]
The individual integrals have this form:
\[
 \int_0^{\frac\pi2}\cos^{2m}\varphi\,d\varphi.
\]
